\documentclass[11pt,a4paper]{article}
\usepackage[T1]{fontenc}
\usepackage[utf8]{inputenc}
\usepackage{graphicx}
\usepackage{booktabs}
\usepackage{amsmath}
\usepackage{hyperref}
\usepackage[margin=1in]{geometry}
\usepackage{natbib}

\hypersetup{colorlinks=true,linkcolor=blue,citecolor=blue,urlcolor=blue}

\title{\textbf{GlycoSMILES2BAP: An Automated Pipeline for Stereochemistry-Preserving Glycan Structure Prediction with AlphaFold 3}}

\author{\textbf{Qiang Xia}\\Zhejiang Xinghe Tea Technology Co., Ltd.\\Hangzhou, Zhejiang, China\\\texttt{xiaqiang@xinghetea.com}}

\date{}

\begin{document}
\maketitle

\begin{abstract}
\noindent\textbf{Motivation:} AlphaFold 3 (AF3) has revolutionized protein structure prediction but standard input formats produce stereochemically incorrect glycan structures. Huang et al.\ (2025) identified that the correct CCD+bondedAtomPairs format requires manual specification taking 30--60 minutes per structure.

\noindent\textbf{Results:} We present GlycoSMILES2BAP, achieving 97.8\% epimer accuracy, 97.4\% anomeric accuracy, and 95.9\% linkage accuracy. Validation against 10 literature errors showed 100\% correction rate. Database-scale processing of 100 structures achieved 94\% success rate at 0.82\,ms per structure.

\noindent\textbf{Conclusions:} GlycoSMILES2BAP eliminates the input format barrier for AF3 glycan modeling.

\noindent\textbf{Availability:} \url{https://github.com/ShawnXiha/glycobap}

\noindent\textbf{Keywords:} AlphaFold 3, glycans, stereochemistry, structure prediction
\end{abstract}

\section{Introduction}

Glycans are essential biological molecules involved in protein folding, cell signaling, and immune recognition. Over 50\% of human proteins undergo glycosylation. GlyTouCan catalogs over 200,000 unique glycan structures.

AlphaFold 3 represents a breakthrough in structure prediction. However, standard input formats fail to preserve glycan stereochemistry, achieving only ${\sim}62\%$ accuracy. The correct CCD+BAP format requires manual atom-by-atom specification.

We present GlycoSMILES2BAP, an automated pipeline converting standard notations to AF3-compatible format.

\section{Methods}

\subsection{Pipeline Architecture}

Four sequential modules: (1) Parser for IUPAC/WURCS conversion, (2) CCD Mapper for 28+ monosaccharide configurations, (3) Anomeric Tracker for C1/C2 handling, (4) BAP Generator for bond specification.

\subsection{Key Features}

Sialic acids use C2 anomeric position; pentoses use O4 ring oxygen. The mapper encodes D/L absolute configuration.

\section{Results}

\subsection{Benchmark Performance}

Achieved 97.8\% epimer accuracy, 97.4\% anomeric accuracy, 95.9\% linkage accuracy on 50 diverse glycan structures.

\subsection{Case Study 3: Error Correction}

\begin{figure}[ht]
\centering
\includegraphics[width=0.9\textwidth]{figures/figure_case_study3.pdf}
\caption{Error correction validation. (A) Distribution of 10 literature error cases. (B) Module contribution analysis. (C) 100\% correction rate across all error types.}
\label{fig:case3}
\end{figure}

Validated against 10 literature-reported errors from PDB structures (Frenz et al., 2018; Jo et al., 2011). All error types corrected: anomeric (4/4), epimer (2/2), linkage (3/3), conformation (1/1).

\subsection{Case Study 4: Database Processing}

\begin{figure}[ht]
\centering
\includegraphics[width=0.9\textwidth]{figures/figure_case_study4.pdf}
\caption{GlyTouCan database processing. (A) 100 structures across 5 categories. (B) 94\% success rate. (C) 1,800$\times$